\begin{textsample}{POS Dim 1 – ac – Score 44.00 – ac/ac\_ef\_5.txt}  \label{ex:f1_pos_176}
1.
REFLEXÕES SOBRE A EDUCAÇÃO FÍSICA A \textbf{introdução} da Educação Física na escola ocorreu, oficialmente, no Brasil, em 1851, com a reforma Couto Ferraz, embora a preocupação com a inclusão de exercícios físicos, na Europa, remonte ao \textbf{século} XVIII, com Guths, J.
J.
Rosseau, Pestalozzi e outros ( BETTI, 1991 ).
Com a reforma realizada por Rui Barbosa, em 1882, houve uma recomendação para que a ginástica fosse obrigatória para ambos os sexos, e que fosse oferecida para as Escolas Normais.
Mas apenas a partir da década de 1920 é que vários estados da federação começaram a realizar suas reformas educacionais e a incluir a Educação Física com o nome mais frequente de ginástica ( BETTI, 1991 ).
A concepção dominante da Educação Física, no início, estava ligada à perspectiva higienista, em que a preocupação \textbf{central} são os hábitos de higiene e saúde, valorizando o desenvolvimento do físico e da moral a partir do exercício.
Já o modelo militarista estava vinculado à formação de uma geração capaz de suportar o combate e atuar na guerra.
A concepção esportista teve seu auge nas Copas do Mundo nas décadas de 58, 62 e 70 e contribuíram para manter o predomínio dos conteúdos esportivos nas \textbf{aulas} de Educação Física.
Segundo Betti ( 1991 ), entre 1969 e 1979 o Brasil observou a ascensão do esporte, a \textbf{razão} de Estado e a inclusão do binômio Educação Física/Esporte na planificação estratégica do governo, muito embora o esporte de alto nível estivesse presente no interior da sociedade desde os anos 20 e 30.
Esse modelo, também chamado de mecanicista, tradicional e tecnicista, foi muito criticado pelos meios acadêmicos, principalmente a partir da década de 1980, embora essa concepção ainda esteja bastante presente na sociedade e na escola nos dias \textbf{atuais}.
As concepções educacionais da Educação Física vêm se modificando ao longo do tempo, de modo que ainda hoje \textbf{influenciam} a formação profissional e suas práticas pedagógicas.
Apesar de ser \textbf{permeada} pelas crises paradigmáticas mencionadas acima, a área vem passando por mudanças significativas que colaboram com a ação docente.
Embora composta de várias concepções, a Educação Física tem apresentado um grande progresso no que \textbf{diz} respeito à sua identidade e relevância no ambiente escolar, pois, de acordo com a Lei de Diretrizes e Bases da Educação n.º 9.394/1996, a Educação Física, integrada à proposta pedagógica da escola, é componente curricular obrigatório de toda a Educação Básica.
A referida lei ainda estabelece, enquanto responsabilidade dos governos Federal, Estaduais e Municipais, a elaboração de novas diretrizes e a definição de conteúdos com base na cientificidade e nas questões do mundo \textbf{contemporâneo}, de modo que, dentre os temas propostos numa perspectiva de inclusão social, estão as diversidades e problemáticas sociais, por meio de uma Base Nacional Comum Curricular, a ser complementada, em cada sistema de ensino e em cada estabelecimento escolar, por uma parte diversificada, \textbf{exigida} pelas características regionais e locais da sociedade, da cultura, da economia e dos educandos ( BRASIL, 2017 ).
Esses vínculos foram determinantes, \textbf{tanto} no que \textbf{diz} respeito à concepção da \textbf{disciplina} e suas finalidades quanto ao campo de atuação e a forma de ser ensinada no contexto \textbf{atual}.
O fato é que, quando se conhecem os pressupostos pedagógicos, é possível melhorar a coerência entre o que se pensa estar fazendo e o que realmente está sendo realizado.
Assim, a área de Educação Física hoje contempla múltiplos conhecimentos produzidos e usufruídos pela sociedade a respeito do corpo e do movimento.
Entre eles, são considerados fundamentais as atividades culturais de movimento, com finalidades de lazer, expressão de sentimentos, afetos e emoções, com possibilidades de promoção, recuperação e manutenção da saúde.
Atualmente, coexistem na Educação Física várias concepções, todas elas tendo em comum a tentativa de romper com o modelo mecanicista, esportivista e tradicional.
São elas : humanista, fenomenológica, psicomotricidade, \textbf{baseada} nos jogos cooperativos, cultural, desenvolvimentista, interacionista-construtivista, crítico-superadora, sistêmica, crítico-emancipatória, saúde \textbf{renovada}, \textbf{baseada} nos \textbf{Parâmetros} Curriculares Nacionais ( PCNs/Brasil, 1998 ), além de outras.
Portanto, é \textbf{tarefa} da Educação Física escolar garantir o acesso dos \textbf{alunos} às práticas da cultura corporal, contribuindo para a construção de um estilo pessoal de exercê -las e oferecendo -lhes instrumentos para que sejam capazes de apreciá -las criticamente.
Compreender a Educação Física a partir de um contexto mais amplo significa entendê -la na sua \textbf{totalidade}, ou seja, compreender que exerce \textbf{influência} e também é \textbf{influenciada} pelas interações que se estabelecem por meio das relações sociais, culturais, políticas, \textbf{econômicas}, religiosas, ético-raciais, de orientação sexual, de geração, de condição física e mental, entre outras, enfatizando o respeito à \textbf{pluralidade} de ideias e a diversidade humana.
Diante disso, a ação pedagógica da Educação Física deve estimular o acesso e reflexão ao acervo de formas e representações do mundo que o ser humano tem produzido, exteriorizadas pela expressão corporal por meio de jogos, brincadeiras, danças, lutas, ginástica, esportes, práticas de aventura, dentre outras, levando em consideração o contexto sociocultural da comunidade educativa ( COLETIVO DE \textbf{AUTORES}, 2012 ). 2.
CONCEITOS-CHAVE E \textbf{ABORDAGEM} \textbf{METODOLÓGICA} O acesso à cultura corporal representa uma oportunidade de diálogo com o ambiente cultural, com tradições, sentidos e significados que as diferentes práticas transformam e ressignificam a cada tempo e contexto.
Aprender dessa cultura significa dialogar com o outro a partir de jogos, brincadeiras, lutas, danças, ginásticas e esportes.
Devemos também considerar que esse diálogo se dará a partir de conteúdos historicamente \textbf{relevantes} da Educação Física e que são portadores de valores.
Alguns deles, como o esporte, muitas vezes são estabelecidos e transmitidos por uma cultura dominante e \textbf{hegemônica}, pela mídia atrelada à indústria do consumo, mantendo o que está previamente estabelecido sem uma reflexão \textbf{crítica}.
Assim, por trás da escolha de cada conteúdo, faz -se presente uma opção política, ética e estética ; por trás da concepção e da metodologia de ensino, revela -se uma perspectiva mais ou menos \textbf{crítica}, mais ou menos emancipatória.
Dessa forma, mais do que o conteúdo tradicional entendido como conjunto de técnicas a serem aprendidas — por exemplo, chutar uma bola, fazer uma estrela, uma determinada técnica de ginástica, ou seja, tudo aquilo que é preciso ensinar explicitamente —, é importante que o professor crie condições para que o \textbf{aluno} aprenda, estabeleça um diálogo entre o jovem e o conhecimento, entre os jovens, entre eles e os professores e com a comunidade e a cultura em que estão todos inseridos.
Aponta -se aqui uma perspectiva \textbf{metodológica} que tem como eixo o diálogo com o conhecimento e não a aprendizagem de um conteúdo estático, acabado e pré-determinado.
É nesta \textbf{capacidade} comunicativa e \textbf{dialógica} que enxergamos a contribuição da Educação Física na formação dos \textbf{alunos} : a possibilidade de um \textbf{aluno} que se quer autônomo e \textbf{crítico}, participativo e emancipado na construção do ambiente de convívio social, a partir das práticas da cultura corporal.
Diante destes preceitos, esta orientação curricular está organizada através de eixos estruturantes de modo a garantir ao \textbf{aluno} o direito de aprender sobre as práticas de cultura corporal, orientando metodologicamente o trabalho pedagógico do professor.
Os eixos temáticos foram pensados a partir de uma fundamentação \textbf{teórica} da Educação Física conhecida por crítico-superadora, descrita e aprofundada no livro “ Coletivo de Autores - Metodologia do Ensino de Educação Física ”.
Os eixos temáticos são : jogo e brincadeira, esporte, ginástica, dança, luta e práticas corporais de aventura.
O último eixo ( práticas corporais de aventura ) surge da necessidade de adequar o Currículo do Estado ao que determina a Base Nacional Comum Curricular - BNCC, visto que o mesmo não era considerado no último documento curricular do Estado, tampouco na fundamentação \textbf{teórica} referida acima.
Cabe destacar que as categorizações das práticas corporais são sugestivas e não têm pretensões de limitar as possibilidades de exploração do universo da cultura corporal, podendo o professor inserir práticas que não estão contempladas ou citadas diretamente neste documento, principalmente as práticas oriundas das experiências dos \textbf{alunos}.
A \textbf{teoria} em questão contribui na concretização do trabalho pedagógico do professor, visto que a \textbf{disciplina} de Educação Física na escola deve contribuir para a formação do \textbf{aluno} em sua \textbf{totalidade}, não se pautando em um currículo tradicional e fragmentado.
Na perspectiva da reflexão sobre a cultura corporal, a dinâmica curricular, no âmbito da Educação Física, tem características bem diferenciadas (... ).
Busca desenvolver uma reflexão pedagógica sobre o acervo de formas de representação do mundo que o homem tem produzido no \textbf{decorrer} da história, exteriorizadas pela expressão corporal : jogos, danças, lutas, exercícios ginásticos, esporte, malabarismo, contorcionismo, mímica e outros, que podem ser identificados como formas de representação simbólica de realidades vividas pelo homem, historicamente criadas e culturalmente desenvolvidas. ( Coletivos de \textbf{autores}, pg. 26, 1992 ).
Como forma de \textbf{materializar} os conteúdos da cultura corporal, \textbf{sugerimos} o \textbf{método} histórico-crítico de educação, tendo como referência Savianni ( 2005 ), que propõe uma \textbf{abordagem} dos conteúdos a partir dos seguintes momentos pedagógicos : 1.
Prática social inicial – levantamento dos conhecimentos prévios dos \textbf{alunos}, o que eles dominam acerca da realidade em forma de conhecimento sincrético ; 2. \textbf{Problematização} – questionamentos acerca da realidade, relacionando com o conteúdo que os \textbf{alunos} deverão discutir e esclarecimento ao longo do processo de \textbf{ensino-aprendizagem} partindo do senso-comum para o senso-crítico ; 3.
Instrumentalização – momento da \textbf{aula} em que o professor apresentará aos \textbf{alunos} os conteúdos ( \textbf{conceituais}, procedimentais e atitudinais ) de relevância histórica, social e científica, de caráter científico e \textbf{sistematizado}, com propostas de atividades diversificadas que permitam ao \textbf{aluno} apropriar -se do conhecimento em direção aos questionamentos realizados no início da \textbf{aula} ; 4.
Catarse – momento de síntese \textbf{avaliativa} em que os \textbf{alunos} irão constatar a apropriação dos saberes adquiridos ; 5.
Prática social \textbf{final} – nova postura do \textbf{aluno} diante da realidade que o cerca, de posse dos conhecimentos adquiridos sistematicamente.
O \textbf{quadro} \textbf{organizador} curricular contém uma tabela de “ propostas de atividades ” que estão organizadas a partir dessa lógica de \textbf{abordagem} dos conteúdos.
Vale ressaltar que tal \textbf{abordagem} não tem caráter estático e linear, podendo ser alterada a critério do professor, de acordo com a realidade apresentada em cada unidade escolar.
A reorganização da orientação curricular de Educação Física ocorreu com a classificação de todos os objetivos de aprendizagem do documento anterior.
A classificação teve como critério os objetivos de aprendizagem que poderiam ser abordados por todos os eixos estruturantes e em todos os anos do Ensino Fundamental I e II e os objetivos que seriam desenvolvidos especificamente em algum ano ao longo do Ensino Fundamental.
Todos os objetivos de aprendizagem que se enquadravam no primeiro critério tiveram alterações no texto que consistia na inclusão do termo referente ao eixo em questão, bem como nos conteúdos a serem abordados para o alcance dos objetivos de aprendizagem e, por fim, nas propostas de aprendizagem.
Os demais objetivos foram mantidos na íntegra.
Os objetivos de aprendizagem foram organizados também de acordo com os campos de atuação que as práticas corporais podem proporcionar ao \textbf{aluno}, quais sejam : cultural/educacional, lazer e saúde, sendo que esta organização está implícita nos objetivos e não literalmente demarcada no \textbf{quadro} \textbf{organizador} curricular.

% matched lemmas: abordagem, aluno, atual, aula, autor, avaliativo, basear, capacidade, central, conceitual, contemporâneo, crítico, decorrer, dialógico, disciplina, dizer, econômico, ensino-aprendizagem, exigir, final, hegemônico, influenciar, influência, introdução, materializar, metodológico, método, organizador, parâmetro, permear, pluralidade, problematização, quadro, razão, relevante, renovar, sistematizar, sugerir, século, tanto, tarefa, teoria, teórico, totalidade
\end{textsample}
